% GENERATED FILE. Edit canonical lessons, then run npm run generate:books.
% canonical-source-hash: fnv1a64:ec0bc1e81a09d847
% canonical-lessons: FA-C119-lubiya, FA-C119-nokhod, FA-C119-sekke, FA-C119-dastband, FA-C119-keravat

\chapter{Beans, Chickpeas, Coin, Bracelet, Tie}
\label{ch:fa-beans-chickpeas-coin-bracelet-tie}
\hlchaptermodality{119}

\begin{chapteropening}
\textbf{By the end of this chapter:} \emph{I can say beans, chickpeas, a coin, a bracelet and a tie.}

Complete the last lesson of chapter 119: I can say beans, chickpeas, a coin, a bracelet and a tie.
\end{chapteropening}

\section[\texorpdfstring{\fa{لوبیا} (lubiyâ)}{لوبیا (lubiyâ)}]{\fa{لوبیا} (lubiyâ) — beans}

\label{lesson:FA-C119-lubiya}



\begin{quote}
Before the new one: say the Persian for a plum, then the Persian for a squash.
\end{quote}

\subsection*{You'll want to know: \fa{لوبیا}}
\textbf{\fa{لوبیا}} — \emph{lubiyâ} — \textquotedblleft{}beans\textquotedblright{}.

\begin{cousinweb}
\emph{lubiyâ} goes back to Greek \emph{lobia}, \textquotedblleft{}pods, beans\textquotedblright{} (EtymDB); Arabic \emph{lūbiyā} is the same word.
\end{cousinweb}

\begin{grammarlens}[title={What you've built}]
One more word for this chapter.
\end{grammarlens}

\subsection*{Your turn}
\begin{itemize}
\raggedright
  \item \emph{Say it:} \emph{lubiyâ}
  \item \emph{Say it:} \emph{lubiyâ}, once more
  \item \emph{Say it:} say \emph{lubiyâ}
  \item \emph{Recall:} say the Persian for a plum, then the Persian for a squash, then say \emph{lubiyâ} again
\end{itemize}

\begin{tcolorbox}[breakable,colback=teal!4,colframe=teal!35!black,title={Before you move on}]
What does \fa{لوبیا} mean? (\textquotedblleft{}Beans\textquotedblright{}.) Say it once more.
\end{tcolorbox}
\section[\texorpdfstring{\fa{نخود} (nokhod)}{نخود (nokhod)}]{\fa{نخود} (nokhod) — chickpeas}

\label{lesson:FA-C119-nokhod}



\begin{quote}
Before the new one: say the Persian for a squash, then the Persian for beans.
\end{quote}

\subsection*{You'll want to know: \fa{نخود}}
\textbf{\fa{نخود}} — \emph{nokhod} — \textquotedblleft{}chickpeas\textquotedblright{}.

\begin{cousinweb}
\emph{nokhod} is a native Persian word, inherited rather than borrowed: it continues Middle Persian \emph{naxōd}, \textquotedblleft{}chickpea\textquotedblright{} (EtymDB).
\end{cousinweb}

\begin{grammarlens}[title={What you've built}]
One more word for this chapter.
\end{grammarlens}

\subsection*{Your turn}
\begin{itemize}
\raggedright
  \item \emph{Say it:} \emph{nokhod}
  \item \emph{Say it:} \emph{nokhod}, once more
  \item \emph{Say it:} say \emph{nokhod}
  \item \emph{Recall:} say the Persian for a squash, then the Persian for beans, then say \emph{nokhod} again
\end{itemize}

\begin{tcolorbox}[breakable,colback=teal!4,colframe=teal!35!black,title={Before you move on}]
What does \fa{نخود} mean? (\textquotedblleft{}Chickpeas\textquotedblright{}.) Say it once more.
\end{tcolorbox}
\section[\texorpdfstring{\fa{سکه} (sekke)}{سکه (sekke)}]{\fa{سکه} (sekke) — a coin}

\label{lesson:FA-C119-sekke}



\begin{quote}
Before the new one: say the Persian for beans, then the Persian for chickpeas.
\end{quote}

\subsection*{You'll want to know: \fa{سکه}}
\textbf{\fa{سکه}} — \emph{sekke} — \textquotedblleft{}a coin\textquotedblright{}.

\begin{cousinweb}
\emph{sekke} is borrowed from Arabic \emph{sikka}, \textquotedblleft{}die for coins, coin\textquotedblright{}. Its Arabic root is \emph{s-k-k} (Lane).
\end{cousinweb}

\begin{grammarlens}[title={What you've built}]
One more word for this chapter.
\end{grammarlens}

\subsection*{Your turn}
\begin{itemize}
\raggedright
  \item \emph{Say it:} \emph{sekke}
  \item \emph{Say it:} \emph{sekke}, once more
  \item \emph{Say it:} say \emph{sekke}
  \item \emph{Recall:} say the Persian for beans, then the Persian for chickpeas, then say \emph{sekke} again
\end{itemize}

\begin{tcolorbox}[breakable,colback=teal!4,colframe=teal!35!black,title={Before you move on}]
What does \fa{سکه} mean? (\textquotedblleft{}A coin\textquotedblright{}.) Say it once more.
\end{tcolorbox}
\section[\texorpdfstring{\fa{دستبند} (dastband)}{دستبند (dastband)}]{\fa{دستبند} (dastband) — a bracelet}

\label{lesson:FA-C119-dastband}



\begin{quote}
Before the new one: say the Persian for chickpeas, then the Persian for a coin.
\end{quote}

\subsection*{You'll want to know: \fa{دستبند}}
\textbf{\fa{دستبند}} — \emph{dastband} — \textquotedblleft{}a bracelet\textquotedblright{}.

\begin{cousinweb}
\emph{dastband} is \emph{dast}, \textquotedblleft{}hand, wrist\textquotedblright{} (EtymDB), with \emph{band}, \textquotedblleft{}band, tie\textquotedblright{} (EtymDB): a band for the wrist. Both are native Persian.
\end{cousinweb}

\begin{grammarlens}[title={What you've built}]
One more word for this chapter.
\end{grammarlens}

\subsection*{Your turn}
\begin{itemize}
\raggedright
  \item \emph{Say it:} \emph{dastband}
  \item \emph{Say it:} \emph{dastband}, once more
  \item \emph{Say it:} say \emph{dastband}
  \item \emph{Recall:} say the Persian for chickpeas, then the Persian for a coin, then say \emph{dastband} again
\end{itemize}

\begin{tcolorbox}[breakable,colback=teal!4,colframe=teal!35!black,title={Before you move on}]
What does \fa{دستبند} mean? (\textquotedblleft{}A bracelet\textquotedblright{}.) Say it once more.
\end{tcolorbox}
\section[\texorpdfstring{\fa{کراوات} (kerâvât)}{کراوات (kerâvât)}]{\fa{کراوات} (kerâvât) — a tie}

\label{lesson:FA-C119-keravat}



\begin{quote}
Before the new one: say the Persian for a coin, then the Persian for a bracelet.
\end{quote}

\subsection*{You'll want to know: \fa{کراوات}}
\textbf{\fa{کراوات}} — \emph{kerâvât} — \textquotedblleft{}a tie\textquotedblright{}.

\begin{cousinweb}
\emph{kerâvât} is borrowed from French \emph{cravate}, \textquotedblleft{}cravat, tie\textquotedblright{} (EtymDB). The French named the neckcloth after the Croatian soldiers who wore it (EtymDB).
\end{cousinweb}

\begin{grammarlens}[title={What you've built}]
One more word for this chapter.
\end{grammarlens}

\subsection*{Your turn}
\begin{itemize}
\raggedright
  \item \emph{Say it:} \emph{kerâvât}
  \item \emph{Say it:} \emph{kerâvât}, once more
  \item \emph{Say it:} say \emph{kerâvât}
  \item \emph{Recall:} say the Persian for a coin, then the Persian for a bracelet, then say \emph{kerâvât} again
\end{itemize}

\begin{tcolorbox}[breakable,colback=teal!4,colframe=teal!35!black,title={Before you move on}]
What does \fa{کراوات} mean? (\textquotedblleft{}A tie\textquotedblright{}.) Say it once more.
\end{tcolorbox}
